% =====================================================================
%  Harald Høffding, Herbert Spencers Filosofi
%  In: Det nittende Aarhundrede. Maanedsskrift for Literatur og Kritik, udgivet af Georg Brandes og Edvard Brandes.  [Aargang 1875] April--September.  København: den Gyldendalske Boghandel (F. Hegel); J. Jørgensen & Co.  Maj 1875, pp. 106-130.
%  Antiqua.
%
%  TRANSCRIPTION -- GENERATED; DO NOT EDIT BY HAND.  Høffding's four pieces of 1875-76 (Herbert Spencers Filosofi; the
%  review of Wilkens; Montaignes Betydning i Etikens Historie; Benedikt Spinoza) are built by ONE pipeline
%  (ebook-method/hoeffding-1875-76/) into a combined working file, which split.py cuts into four transcriptions.
%  Rebuild from the repo root:   sh ebook-method/hoeffding-1875-76/rebuild.sh
%
%  STATUS: COMPLETE (2026-10-15): the volume's title page and the whole text, pp. 106-130, with its 13 notes.
%
%  EDITION.  The only printing.  Text of record: the Internet Archive's own Scribe scan of the University of Illinois copy, detnittendeaarhu1875unse (500 ppi colour JP2); extract ~/bibliotek/Høffding, Harald/1875-nittende-aarhundrede-apr-sep-pp106-130.pdf (SCANS.tsv, sha256), made by ia_extract.py
%  (title leaf, Indhold leaves, the article).  SECOND COPY (voter only): Google's scan of the Harvard copy, nittendeaarhund05unkngoog (its ABBYY OCR).
%  Page range from the volume's Indhold and the scan.  The note on the title (a lecture in Studenterforeningen, 13 February 1875) is the print's.
%
%  ACCURACY (2026-10-15; one collation for all four texts, 57 working pages, ~17,100 words).
%  (a) The DRAFT is our own tesseract reading (tessdata_best: dan for Det nittende Aarhundrede, script/Fraktur for Nær og
%      Fjern; 300 dpi) of the Archive's JP2 images, laid over the working PDF and turned into LaTeX by draft.py.  It was
%      chosen over the Archive's own OCR (tesseract 5.2/5.3, hOCR) by measurement against the Harvard copy's ABBYY OCR
%      (words unmatched: 3.0% against 9.5% in the Montaigne volume; the Archive reads aa as åa and long s as f).  The
%      Archive's OCR is the WITNESS: collated word by word (collate/check.py): 900 disagreements, 76 settled by check.py.
%  (b) Rules (rules.py: the Archive's å for a, long s read f, glyph faults against attested words, 326) and a THIRD reading
%      voted: the ABBYY OCR of ANOTHER COPY (Google's scan of the Harvard volume), by context and by word alignment
%      (rules.py harv-*, vote2.py hv-*), never on figures and never making a word the corpus does not attest.  AUDIT: 30
%      letter-changing rule/vote decisions read blind at the image: 0 contradicted (6 undetermined crops).
%  (c) The rest of Det nittende Aarhundrede read at the image BLIND (round b1: 210 crops on 12 sheets, five readers, 23/24
%      planted controls right; ~80 decided by hand from the readers' lines).  OPEN 0.
%  (d) Every FOOTNOTE read whole at the image (readers N1, N2: 26 notes on 23 pages), set from the reading; the note marks
%      checked at zooms (four the draft had misplaced, moved).  STRUCTURE read at the image (reader K): the title pages, the
%      heads, the endings and signatures, the verse on p. 116.
%  (e) Nær og Fjern (FRAKTUR, two columns, no second copy): each column PROOFREAD WHOLE at the image against the draft by
%      two readers (P1, P2; 3 of 4 planted errors found): 171 corrections (proofapply.py); then an unattested-word sweep
%      and a long-s / ø sweep, every doubtful word read at a zoom (proof-fixes.tsv, final-fixes.tsv).
%  (f) LETTERSPACING (\emph{} here) MEASURED at the image (spacemeasure.py, white runs between letters, threshold 5 px: 31
%      runs), every run judged at the image by reader K (13 kept in the body; the notes' spacing read by N1/N2).  Italics:
%      none found by stroke slant (check.py --italics).
%  (g) QUOTATION MARKS: the print quotes „...“ only (every quotation read at the image in b1, the notes and the structure);
%      both OCRs read „ as , or ,, and “ as ' ” " »: set by class (quotes.py), outside the notes.
%  (h) Sweeps on the collated text, every doubtful spot read at a zoom: unattested words (229 words in no other
%      transcription: 12 words mended, the misprints below confirmed); specks both OCRs had read as marks (specksweep.py,
%      marks the Harvard copy lacks mid-sentence: 16 removed).  The TRANSLATORS' odd spots (64) were read at zooms before
%      the translation was finished: 22 more faults mended (specks, ɔ: read as 2:, a running head and a footnote left in
%      the text, two paragraph breaks at page turns), 3 misprints confirmed (final-fixes.tsv).
%  (i) SAMPLED PAGE CHECK (2026-10-15, budget-limited to 3 pages): pp. 113 (Spencer), 217 (Montaigne) and one Nær og Fjern
%      column (p. 2, right), chosen at random (seed 20261015, one per stratum), transcribed BLIND at the image by one
%      reader (1,143 words) and compared word by word (pagecheck.py): 10 differences; decided at zooms: 3 reader slips;
%      7 errors of the collated text (2 words: „for“ for „før“, „skikkelig“ for the print's „stikkelig“; 5 marks: four
%      specks both OCRs had read as stops or quotes, one stop hidden by a planted control) = ~6 per 1000 words, all mended.
%      The specks led to specksweep.py over the whole of Det nittende Aarhundrede: 16 more found at zooms and removed.
%  PRINTED AS IS (misprints kept, marked at the line):
%    p. 106: „filososofisk“ for filosofisk; sweep zoom
%    p. 108: „Menesket“ for Mennesket; sweep zoom
%    p. 111: „Fænomernes“ for Fænomenernes; sweep zoom
%    p. 111: „Grundforudsætnirg“ for Grundforudsætning; sweep zoom
%    p. 129: „Tisrum“ for Tidsrum; reader b1-04-01
%  CONVENTIONS.  Page turns: % --- p. N --- and \opage{N} (the printed page); a word divided over a page turn is
%  split with %.  Footnotes as \footnote at their marks, numbered per page as printed (*), **)).  \emph{} = letterspacing.
%  The running heads and the sheet signatures are not transcribed.  The ø rule (CLAUDE.md).
% =====================================================================
\documentclass[12pt,a4paper,openany]{book}
\usepackage[utf8]{inputenc}
\usepackage[T1]{fontenc}
\usepackage{libertinus}
\usepackage{libertinust1math}
\usepackage[danish]{babel}
\usepackage[protrusion=true,expansion=true]{microtype}
\usepackage{setspace}
\usepackage[margin=1.3in]{geometry}
\usepackage{fancyhdr}
\usepackage{hyperref}
\hypersetup{pdftitle={Herbert Spencers Filosofi}, pdfauthor={Harald Høffding}, hidelinks}
\pagestyle{fancy}
\fancyhf{}
\fancyhead[LE,RO]{\thepage}
\fancyhead[C]{\textit{H. Høffding: Herbert Spencers Filosofi (1875)}}
\fancypagestyle{plain}{\fancyhf{}\fancyhead[LE,RO]{\thepage}}
\setstretch{1.3}
\usepackage{marginnote}
\usepackage{graphicx}
\DeclareUnicodeCharacter{0254}{\reflectbox{c}}   % ɔ (id est)
\makeatletter
\renewcommand{\thefootnote}{\ifcase\value{footnote}\or *\or **\or ***\else\arabic{footnote}\fi)}   % the print marks each page's notes *), **)
\makeatother
\newcommand{\opage}[1]{\marginnote{\footnotesize\textit{[#1]}}}
\newcommand{\signature}[1]{\par\smallskip\hfill #1\hspace*{2em}\par}
\pdfstringdefDisableCommands{\renewcommand{\footnote}[2][]{}\renewcommand{\opage}[1]{}}

\begin{document}

\begin{titlepage}
\begin{center}
{\LARGE\bfseries Det nittende Aarhundrede.}\\[1em]
{\small\bfseries \emph{Maanedsskrift}}\\[0.3em]
{\small for}\\[0.3em]
{\large\sffamily Literatur og Kritik.}\\[1.2em]
{\small Udgivet af}\\[0.3em]
{\large Georg Brandes {\small og} Edvard Brandes.}\\[0.6em]
\rule{3em}{0.4pt}\\[0.6em]
{\small\bfseries 1875. April---September.}\\
\vfill
% [printer's device: a ring with the monogram GB]
{\large\sffamily \emph{København}.}\\[0.4em]
{\small Forlagt af den Gyldendalske Boghandel (F. Hegel).}\\
{\small Trykt hos J. Jørgensen \& Co.}\\[0.4em]
{\small 1875.}
\end{center}
\end{titlepage}
\clearpage

% --- p. 106 ---
\setcounter{footnote}{0}%
\opage{106}\begin{center}{\Large\sffamily\bfseries Herbert Spencers Filosofi.\footnote[1]{I det Væsentlige et Foredrag, holdt i Studenterforeningen den 13de Februar 1875.}}\\[0.5em]\rule{4em}{0.4pt}\end{center}\medskip

Naar man betegner en videnskabelig Bestræbelse
som filososofisk, vil det kunne være to tilsyneladende % PRINTED AS IS (filososofisk: for filosofisk; sweep zoom)
meget forskellige Ejendommeligheder, der berettige den
til dette Navn. Det vil enten være, fordi den har en
afgjort kritisk Karakter, vender sig tvivlende og prøvende
mod hidtil faststaaende Forudsætninger og Resultater,
eller fordi den søger at sammenfatte alle Erkendelsens
Resultater og Synspunkter i en Grundanskuelse, altsaa
formedelst dens systematiske Karakter. Dog udelukke
disse Egenskaber ikke hinanden; de findes til en vis
Grad forenede hos enhver Tænker, kun at alt efter hans
Karakter og efter den Tidsalder, i hvilken han lever,
enten Kritiken eller Systemet vil være det Overvejende,
--- og Filosofiens ideale Fuldendelse vilde være en gennemført
Kritik og et Alt omfattende System. Det vil ogsaa
være let at vise, at Kritiken og Systemet gensidig vise
hen til hinanden. Kritik er Sondring, Analyse. Den
kan da ikke hvile, førend den har fundet alle Elementer,
alle Sider i de Genstandes Væsen, med hvilke den beskæftiger
sig. Men Tanken kan ikke heller slaa sig til
Ro ved det Adskilte og Usammenhængende; Kritikens
og Analysens Resultater bruges derfor atter som Stene
% --- p. 107 ---
\setcounter{footnote}{0}%
\opage{107}til en ny Bygning, som er des fastere og sikrere, jo dybere
og grundigere den foregaaende Kritik har været. Disse
to forskellige Retninger ere for Tænkningen, hvad Indaanding
og Udaanding er for Organismen, begge lige
uundværlige, og paa deres rytmiske Skiften beroer Tankens
Sundhed.

I Almindelighed vil man sikkert mene, at i vor
Tid gaaer Bevægelsen overvejende i den kritiske Retning.
Vor Tid, usikker som den er paa sig selv i Theori som
i Praksis, Tvivlens og Pessimismens Periode, synes ikke
gunstig for Systemer og Syntheser. Og dog turde der,
i det Skjulte, ved Bidrag fra mange forskellige Sider
være samlet Stene til en Bygning, som engang --- hvorledes
det saa end gaar med de foreløbige Udkast --- vil
kunne rejse sig som et mægtigt Tempel, et Udtryk for
alle Menneskeaandens ideelle Bestræbelser.

Den, der sammenligner det 19de Aarhundrede med
dets Forgænger, vil ikke kunne negte, at de forholde
sig til hinanden som System til Kritik. Det 18de Aarhundredes
Historie er en Opløsningsproces; det 19de
stræber under mange forskellige Former efter Organisation.
Hvad særlig Filosofien angaar, saa have de vigtigste Forsøg
i det 18de Aarhundrede --- Locke, Condillac og Kant ---
en afgjort kritisk og analytisk Karakter, medens den
systematiske Tendens i vort Aarhundrede ligeledes repræsenteres
af de mest fremragende Tænkere hos de tre
Hovednationer: Hegel, Comte og Spencer. Vi finde her
en og samme Grundopgave behandlet paa tre forskellige
Maader.

For Hegel, der stiller sig paa den aprioriske Spekulations
Standpunkt, er det Tankens Liv, der pulserer i
Alt. De Stadier, Tanken gennemgaar for at naa sin
Fuldendelse i Begrebet, ere Et med dem, gennem hvilke
al Virkelighed udfolder sig: de to Bevægelser falde
sammen. Skønt Hegel gaar ud fra det rent Logiske,
har han dog grebet det frugtbare Synspunkt at se dette
% --- p. 108 ---
\setcounter{footnote}{0}%
\opage{108}som Bevægelse, Proces, Udvikling; derved blev det ham
muligt at forsøge en Systematisering af Erkendelsens
hele Indhold. Men selv om han havde Ret i sin Opfattelse
af Tankeprocessen, havde han dog ikke Ret i
uden videre at overføre de for denne Proces gældende
Love paa det Reale.

Comte søgte at frigøre sig for al Spekulation og
stillede sig paa den positive Videnskabs Standpunkt.
Han vilde forklare Menesket ved Verden, ikke Verden % PRINTED AS IS (Menesket: for Mennesket; sweep zoom)
ved Mennesket. Og dog er han beslægtet med Hegel
ved sin afgjorte systematiske Tendens, sin Higen efter
og Sans for Totalitet\footnote[1]{Hegel og Comte ytrede gensidig Sympati for hinandens Bestræbelser, rigtignok uden udtømmende Kendskab til disse. Se Comtes Brevveksling med Eichthal hos Littré: Auguste Comte et la philosophie positive.}. Han vilde ophæve de enkelte
Videnskabers Isolation ved i en almindelig Klassifikation
at samle deres Principer og Resultater for at vise, hvorledes
de tilsidst gaa op i den ene, Alt omfattende Videnskab:
Menneskehedens Videnskab. Hans Fejl bliver den
modsatte af Hegels: hos denne blev Virkelighedens Mangfoldighed
voldsomt presset ind under Tankens Enhed; hos
Comte taber Tankens Enhed sig i Virkelighedens Mangfoldighed,
thi den blotte Klassifikation kan ikke tilvejebringe
en virkelig Enhed.

Spencers Forsøg, som vi her nærmere skulle beskæftige
os med, staar paa en Maade midt imellem Hegels
og Comtes. Filosofi, siger han, er den til Enhed og
Helhed udviklede Erkendelse. Vi have altsaa ikke naaet
Maalet, saa længe de enkelte Videnskaber og deres Indhold
endnu, som hos Comte, danne en Rækkefølge; vi
maa finde én fælles Lov, ét Princip, hvori Alt, hvad vi
kende, finder sit Udtryk. Men dette Princip finder
Spencer ikke, som Hegel, umiddelbart i Tankens eget
Væsen. Han finder det ad Induktionens Vej, ved Sammen%
% --- p. 109 ---
\setcounter{footnote}{0}%
\opage{109}ligning af Fænomenerne paa alle forskellige Omraader,
lige som Kepler i sin Tid fandt sine berømte Love ved
Iagttagelse af Planeten Mars paa forskellige Punkter af
dens Bane. Men hvad der saaledes viser sig at være en
almindelig Lov, gældende for alle Fænomener, maa saa
ogsaa vise sig at staa i den inderligste Sammenhæng med
de sidste Principer for vor Erkendelse i det Hele.

Vi forstaa først et Fænomen, naar vi kunne følge
det fra dets første Begyndelse og paa de forskellige Trin
af dets Udvikling. Lige som Kendskab til et Folk vil sige
Kendskab til dets Historie, Kendskab til et Menneske,
Kendskab til hans Familie, Barndom, Ungdom og Manddom,
saaledes kende vi ogsaa først et Træ eller et Dyr, naar vi
vide, hvorledes de efterhaanden udfolde sig fra de
usynlige Spirer, og et kemisk Stof, naar vi vide, af hvilke
Elementer det er sammensat. Nu er det klart, at naar
et Fænomen opstaar, vil dette sige, at Elementer, som
før eksisterede i adspredt Tilstand samles og kombineres
paa en ejendommelig Maade. Vandet opstaar ved en vis
bestemt Kombination af Ilt- og Brintatomer, Planten og
Dyret ved organisk Assimilation af visse kemiske Stoffer,
et Menneskes Karakter ved Vekselvirkning mellem hans
medfødte og nedarvede Natur med de ydre Forhold o. s. v.
Ethvert Fænomen opstaar altsaa ved en Koncentrationsproces.
De enkelte Elementer miste deres selvstændige
og uafhængige Bevægelse og bringes for en Tid i Ligevægt
paa en vis bestemt Maade. Men ogsaa kun for en
Tid. Thi som ethvert Fænomen gennemgaar en Udviklingsproces,
saaledes gennemgaar det ogsaa en Opløsning.
Koncentrationen bringer Delene i Ligevægt; men denne
varer ikke ved, den afløses af en Adsplittelse af de Elementer,
der for en Tid havde kombineret sig til Dannelse
af et ejendommeligt Hele. Ethvert Fænomens Historie
bestaar i en Udviklingsproces og en Opløsningsproces.

Dog kunne vi endnu noget nøjere karakterisere disse
Processer. Naar Udviklingen er en Koncentration af
% --- p. 110 ---
\setcounter{footnote}{0}%
\opage{110}forhen spredte Elementer, saa er det for det Første
klart, at den betegner en Overgang fra en mindre sammenhængende
til en mere sammenhængende Tilstand. Ifølge
den Kant-Laplaceske Hypothese har saaledes vort Solsystem
udviklet sig ved, at den oprindelige Taagemasse
sammentrængtes til Dannelse af enkelte Kloder. Jordklodens
Udvikling bestaar i en Fortætning. Det organiske
Liv beror paa den nøje Sammenhæng mellem de Elementer,
der ved Ernæringen optages fra Omverdenen.
Det menneskelige Samfund er paa sit laveste Trin en
Flok Vilde, der ikke blot mangler indre Sammenhæng,
men ogsaa er isoleret fra andre Samfund; jo videre Udviklingen
skrider frem, des mere Sammenhæng indenfor
det enkelte Folk og mellem alle enkelte Folk indbyrdes.

Det, der saaledes under Udviklingens Gang bringes
i nøjere Sammenhæng, gaar for det Andet over fra en
ensartet til en uensartet Tilstand. I den oprindelige Taagemasse
var der ingen væsentlige Forskelligheder; disse opstaa
først ved de enkelte Kloders Koncentration. Under
Jordklodens Udvikling danner sig først en Modsætning
mellem Skorpe og Kærne, saa paa selve Skorpen en Modsætning
mellem Bjerge og Sletter, Have og Fastlande.
Under Organismens saa vel som under Samfundets Udvikling
indtræder der en Arbejdets Deling, en gennemgaaende
Specialisering af Opgaver og Funktioner.

For det Tredje bestaar Udviklingen i en Overgang
fra det Ubestemte til det Bestemte. Orden afløser overalt
Kaos; udprægede individuelle Former danne sig:
saaledes de forskellige Kloder, de forskellige Nationer,
Stænder og Personligheder.

Disse tre Love, forskellige Former for den almindelige
Udviklingslov, ere de Hovedsynspunkter, hvortil enhver
Betragtning af Fænomenerne fører os. De ere de
almindelige Generalisationer af al vor Erfaring. Men
det Spørgsmaal rejser sig nu, om ikke en dybere Begrundelse
og Afledning af disse Love skulde være mulig,
% --- p. 111 ---
\setcounter{footnote}{0}%
\opage{111}lige som de keplerske Love jo først fik deres rette Begrundelse
ved, at Newton førte dem tilbage til Tyngdeloven.

Ethvert Fænomen henføre vi nødvendigvis til et eller
andet tidligere Fænomen; heri bestaar, som vi have set,
den Forklaring, vi kunne opnaa af Fænomenerne; Forklaring
af et Fænomen er Paavisning af dets Udvikling.
Det ene Led i Fænomernes uendelige Række belyser det % PRINTED AS IS (Fænomernes: for Fænomenernes; sweep zoom)
andet. Bevidsthedens Grundvirksomhed er en bestandig
Overgang fra en Fornemmelse eller Forestilling til en
anden. Men saa er det da ogsaa klart, at vor Bevidsthed
kun kan fungere, eller med andre Ord, at vi kun kunne
tænke, saa længe en saadan Overgang fra et Led til et
andet er mulig. Et Fænomen af den Beskaffenhed, at
det i sin Fremtræden og i sit hele Væsen aldeles ikke
hænger sammen eller viser hen til nogetsomhelst andet,
er en Umulighed, vi ikke engang kunne konstruere for
vor Fantasi. Tænkningens Grundforudsætnirg er derfor % PRINTED AS IS (Grundforudsætnirg: for Grundforudsætning; sweep zoom)
den, at der ingen absolut Begyndelse eller Ende er, ---
ingen Overgang fra Intet til Noget eller fra Noget til
Intet: en saadan Overgang vilde jo være selvmodsigende,
vilde være et Forhold med kun ét Led!

Er en absolut Opstaaen og Forgaaen umulig, saa
maa Materiens Sum ogsaa stedse være den samme, og
 Bevægelsen maa være kontinuerlig. Materien kende vi
nu kun gennem den Modstand, den gør imod vor Bevægelse,
og Bevægelsen kun ved den Forandring, den bevirker
i vor Bevidsthed: Kraftens Begreb ligger saaledes
til Grund for vore Begreber om Materie og Bevægelse,
og Kraftens Bestaaen er altsaa den Grundforudsætning,
hvorpaa al Erkendelse hviler. Den Sætning, at Kraften
bestaar, er den eneste, der gaar ud over al Erfaring,
fordi den ligger til Grund for al Erfaring.

Den moderne Kemi og Fysik have vel bevist, at de
materielle Elementer aldrig blive til Intet, men ved enhver
Forsvinden af et Legeme blot gaa over i nye Kombinationer,
og at ingen Kraft bliver til Intet, men ved
% --- p. 112 ---
\setcounter{footnote}{0}%
\opage{112}sin tilsyneladende Forsvinden kun omsættes til en anden
Kraft. Men selve Sætningen om Kraftens Bestaaen kan
ikke bevises ad induktiv Vej af den simple Grund, at
enhver Induktion forudsætter den. Ethvert kemisk og
fysisk Eksperiment hviler paa denne Sætning: saasnart der
skal vejes eller maales, maa der lægges en Krafténhed
til Grund; men hvis denne Krafténhed kunde forandre
sig, vilde Eksperimentets Resultat være illusorisk. Al
videnskabelig Forsken gaar derfor ud fra denne Forudsætning.

Den dybeste Begrundelse, vi derfor kunne vinde af
nogen empirisk Lov, bestaar i at føre den tilbage til Sætningen
om Kraftens Bestaaen. Lad os da igen se paa
Fænomenernes Udviklingslove. Oprindelig ere Elementerne
spredte i en ensartet Masse. Men da de enkelte
Dele af denne modtage forskellige Paavirkninger fra Omgivelserne,
saa ville de efterhaanden ogsaa selv blive
forskellige. Dersom nemlig forskellige uden fra virkende
Kræfter havde samme Indflydelse paa de ensartede Elementer,
saa vilde dette stride mod Kraftens Bestaaen:
Kræfternes Forskellighed vilde ikke lægge sig for Dagen
i en Forskellighed i Virkningen. Og naar der
nu altsaa danner sig Forskelligheder indenfor den oprindelig
ensartede Masse, saa vil paa den anden Side
samme Kraft have forskellige Virkninger paa de forskellige
Dele. Disse to Sætninger: at forskellige Kræfter
øve forskellig Indflydelse paa ensartede Dele, og at samme
Kraft øver forskellig Indflydelse paa forskellige Dele, ---
som begge udspringe af Grundsætningen om Kraftens
Bestaaen, --- indeholde Grunden til Udviklingsprocessen
som Overgang fra ubestemt og usammenhængende Ensartethed
til bestemt udpræget og sammenhængende Forskellighed.

Udviklingsprocessen varer nu ikke bestandig. Tilsidst
naas der en Tilstand, hvor der er Ligevægt mellem
% --- p. 113 ---
\setcounter{footnote}{0}%
\opage{113}Fænomenet og dets Omgivelser. De udenfra virkende
Kræfter virke nu ikke længer til Udvikling; de enkelte
Elementer have nu mistet deres selvstændige Bevægelse;
der kan da ikke skrides videre frem i Retning af Koncentration.
Men Ligevægten er kun Begyndelsen til Opløsningen:
thi de ydre Kræfter høre ikke op at virke,
og deres Indflydelse vil nu bestaa i, at Elementerne
spredes og igen gaa over i selvstændig Bevægelse. Ogsaa
dette følger lige frem af Kraftens Bestaaen. Naar Opløsningen
da har naaet sin Fuldendelse, ville ydre Kræfter
igen kunne virke koncentrerende --- og saaledes i det
Uendelige.

Universets Liv i alle dets Dele og Former er saaledes
en uendelig Rækkefølge af Udvikling og Opløsning.
Det er den mest almindelige Lov, hvortil Erfaringen kan
føre os, og den finder, som vi have set, sin Begrundelse
i selve vor Bevidstheds Grundlov.

Men Spencer fastholder paa det Bestemteste, at vi
herved ikke naa ud over det Fænomenale og Relative.
Det Absolute, det i og for sig Værende kunne vi hverken
forestille eller tænke: vor Bevidstheds Virksomhed
er jo en Overgang fra et Led til et andet --- men det
Absolute, det i sig Fuldendte og Afsluttede, tilsteder
ingen saadan Overgang, da det ikke har Noget udenfor
sig. Alligevel nødes vi til at antage et absolut Princip:
ellers vilde selve det Relative og Betingede jo blive absolut.
Der maa eksistere Andet end blotte Skranker og
Betingelser. Vi have desuden set, at der under al Fænomenernes
Skiften, deres Oprindelse og Opløsning, dog
bestaar Noget --- Kraften i sig selv, som ikke kan begynde
eller ophøre. Det Absolute, det i sig Eksisterende
er den Kraft, som rører sig i Alt, --- hvis Væsen vi ikke
kende og ikke kunne kende, men som vi maa antage
som sidste Grund og Forudsætning for Alt, hvad vi
kunne kende.

% --- p. 114 ---
\setcounter{footnote}{0}%
\opage{114}Idet Videnskaben saaledes erkender, at den, hvor
vidt den saa naar i Erkendelse af Fænomenernes Sammenhæng,
aldrig vil kunne naa den dybe Kilde, hvoraf de
udvælde, medens vi dog med Nødvendighed føres til at
antage dens Eksistens, bliver en Forening med den
religiøse Bevidsthed mulig. Ogsaa Religionerne hævde,
at Tilværelsens Kærne er et Mysterium for Menneskene,
--- men de hilde sig samtidig i den Modsigelse at paastaa
at være i Besiddelse af særdeles nøje Oplysninger
om dette Mysterium. Fetischdyrkeren har det Absolute
for sig i Træklodsen eller Stenen, som han slaar, naar
den ej vil høre ham; paa højere Stadier er det Dyre- eller
Menneskeskikkelsen, man ifører det; Verdens Historie
er endnu for Mange den absolute Aands Selvbiografi.
Overfor alle positive Religionsformer har Fichte Ret,
naar han siger, at en Gud, vi kunde begribe, ikke vilde
være nogen Gud, men en Afgud. Gennem Religionernes
Historie kan man ogsaa forfølge en fremskridende Bevidsthed
om denne Sandhed. Men endnu den Dag idag
ere de saakaldte Fromme, der klamre sig til de overleverede
endelige Udtryk for det Ubetingede og for at
hævde disse forhærde sig mod Videnskabens Resultater,
hvorved de gøre sig skyldige i den værste Vantro: Frygten
for at Sandheden skulde være skadelig, --- ere disse
saakaldte Fromme mindre religiøse end Forskeren, der
véd, at alle slige Former kun ere Symboler.

Ogsaa paa dette afgørende Punkt indtager Spencer
en Stilling, der ligger midt imellem Hegel og Comte.
Hegel finder i de spekulative Begreber, i hvilke Tilværelsens
Kærne udtømmende erkendes, ogsaa den positive
Religions Sandhed indeholdt. Comte ansaa derimod
det religiøse Standpunkt for fuldstændig overvundet; den
eneste Religion, der nu kunde være Tale om, var ifølge
hans Overbevisning Humanitetens Religion, der bestaar i
den inderlige, endogsaa mystiske Hengivelse til Menneskeheden
og dens Udvikling. Spencer hævder herimod, at
% --- p. 115 ---
\setcounter{footnote}{0}%
\opage{115}den religiøse Følelse stedse vil vende sig mod en absolut
Grund til al Tilværelse, kun at den maa opgive ethvert
Forsøg paa at lære os Noget om denne absolute Grund.
Universet selv i sin Rigdom og i sin uendelige rytmiske
Bevægelse er den eneste Aabenbaring af den evige,
uerkendelige Kraft, som rører sig i Alt.

Uagtet jeg her kun vil fremstille Spencers Tankegang
og ikke min egen Stilling til den, kan jeg dog ikke
holde den Bemærkning tilbage, at Spencer vanskelig vil
kunne holde Udviklingens Begreb ganske udenfor al Anvendelse
paa det Absolute. Naar det er en Grundlov
for Alt, at det gennemgaar Udvikling og Opløsning, saa
kan Tilværelsens Kærne ikke være aldeles fremmed for
denne Bevægelse. Uden at ville begive os ud paa nye
spekulative Eventyr kunne vi dog, og synes vi at maatte
antage, at den rytmiske Bevægelse har Betydning for
selve det absolute Princip. Det er lige saa berettiget at
betegne dette som Udviklingens Princip, som at bruge
Udtrykket Kraft om det; Udviklingens Lov har jo ifølge
Spencer samme Universalitet og Nødvendighed som Kraftens
Begreb. Kun under denne Forudsætning synes ogsaa
den religiøse Følelse at kunne se Universet som virkelig
Aabenbaring af det Ubetingede; og kun under denne
Forudsætning kan Resignationen under Verdenslovene
blive en Kilde til ny Begejstring og Stræben.

Spencer selv har iøvrig givet denne Begejstring Udtryk
i følgende Udtalelse, fremkaldt ved den Modsætning,
Forskeren stedse føler sig staaende i overfor de positive
Religioners Tilhængere, hvormeget han end anerkender
eller maaske netop fordi han anerkender det relativt Berettigede
i deres Tro. „Han [den fremskredne Tænker]
kan, som ethvert andet Menneske, med Rette betragte
sig som et af de utallige Organer, gennem hvilke den
ubekendte Aarsag virker, og naar den ubekendte Aarsag
frembringer en vis Overbevisning hos ham, bemyndiges
han derved til at udtrykke den i Tale og  Handling.
% --- p. 116 ---
\setcounter{footnote}{0}%
\opage{116}Thi --- for at anvende Digterens Ord i deres højeste
Betydning ---
\begin{verse}\small
\hspace*{6em}Naturen kan dog ikke\\
forbedres ved et Middel, som Naturen\\
ej selv har skabt; og over denne Kunst,\\
der, som du siger, hjælper paa Naturen,\\
er der en Kunst, Naturen selv har skabt.
\end{verse}
\begin{flushright}\small\textbf{(Shakspeare: Vintereventyr.)}\end{flushright}

Derfor vil den Vise ikke betragte den Overbevisning,
der rører sig i ham, som tilfældig. Han vil uden Frygt
udtale den højeste Sandhed, han kan se, fordi han véd,
at hvad saa Følgen bliver, spiller han derved den rette
Rolle i Verden, --- fordi han véd, at \emph{kan} han bevirke
den Forvandling, han stræber efter, godt! --- kan han
\emph{ikke}, ogsaa godt, skønt ikke \emph{saa} godt! (First principles
3d ed. p. 123).“ Her antyder Spencer aabenbart et andet
Forhold mellem det absolute Princip og Udviklingens
Proces end i de ovenfor givne Bestemmelser.

Vi vende tilbage til Udviklingens Love. Det er paa det organiske og det sociale Omraade, at den mest iøjnefaldende
lægger sig for Dagen. Udviklingsperioderne ere
her kortere, Betingelserne lettere overskuelige. Det er
her, at Udviklingens Ide først opgik for Tanken, navnlig
ved at man blev opmærksom paa Videnskabernes og
Samfundstilstandenes Udvikling gennem Tiderne og paa
den enkelte Organismes Vækst. I sit første betydeligere
Skrift „Social Statics“ (1850) udtalte Spencer allerede,
at den samme Lov gjaldt for den organiske og den sociale
Udvikling. Det er især den moderne Embryologi, grundlagt
af Harvey, Wolff og v. Baer, som har givet ham
Impulsen til en af hans vigtigste Ideer: Udviklingen som
Overgang fra en ensartet til en forskelligartet Tilstand.
Skridt for Skridt førtes han saa til at anvende denne
Lov paa Universet i det Hele.

Spencer var saaledes fuldt beredt til at gaa ind paa
den storartede Hypothese, som Darwin opstillede i sit
% --- p. 117 ---
\setcounter{footnote}{0}%
\opage{117}berømte Værk om Arternes Oprindelse (1859). Darwin
nævner ham ogsaa her som en af sine Forgængere. Men
Spencer lægger i dette Spørgsmaal ikke Hovedvægten
paa den samme Faktor, som Darwin. Denne betragter
de særegne Arter som opstaaede under den gennem utallige
Generationer fortsatte Kamp for Livet derved, at i
enhver Slægt de Individer, der ikke egne sig til og ikke
have Kræfter til at hævde deres Eksistents overfor Omgivelserne
bukke under, fortrænges af de stærkere eller
mere skikkede og udelukkes fra Forplantning. Da Slægten
saaledes bestandig forplanter sig gennem de til at
leve under de givne Forhold mest skikkede Individer,
vil der efterhaanden udvikles en Form, som er i saa
god Harmoni med Betingelserne som muligt. Dette er,
hvad Darwin kalder den naturlige Udvælgelse eller Kvalitetsvalget.
Ifølge Spencer har Darwin lagt for megen
Vægt paa denne indirekte Udvikling. Selv tager han sit
Udgangspunkt fra den Erfaring, at Øvelsen, Funktionen
udvikler Organet. Naar en Muskel hyppig bruges, behøver
den mere Næring; der strømmer da mere Blod til
den end før, og den vil derved tage til i Omfang. Nu
er der vel i det enkelte Individ en bestemt Grænse for
denne Tiltagen; da kommer et Punkt, hvor Blodkarrene
ikke ere store nok til at yde den udkrævede forøgede
Tilførsel. Men det er ogsaa en Kendsgerning, at en ved
forøget Funktion udviklet Struktur kan gaa i Arv til
Efterkommerne. Den ene Slægt fortsætter da, hvad den
foregaaende har begyndt, og saaledes kan Organisationen
Skridt for Skridt udvikles til en saadan Form, som de
ydre Betingelser kræve. Funktion kan \emph{forsaavidt} siges
at gaa forud for Struktur. De laveste Organismer f. Eks.
Rhizopoderne, ere saaledes Eksempler paa Liv og Fungeren
uden egentlig Organisation\footnote[1]{Det er naturligvis ikke Spencers Mening, at Funktionen kan være uafhængig af Strukturen. I saa Fald vilde G. H. \emph{Leves} \opage{118}have Ret i sin Polemik mod den Sætning, at Funktion gaar forud for Struktur. (Fortnigthtly Review. New Serie, vol. IV p. 69---73). Hvad Spencer mener er sikkert kun det, at Funktion kan blive Aarsag til Udvikling af nye Strukturforhold, hvilket lige frem stemmer med Erfaringen. Se hans „Principles of biology“, I p. 185---199.}% the note runs over from p. 1117 to p. 1118 (joined by draft.py)
.

% --- p. 118 ---
\setcounter{footnote}{0}%
\opage{118}Spencer indrømmer, at den indirekte Udviklingsproces,
som bestaar i de ufuldkomne Individers Undergang og
Slægtens Forplantning gennem de fuldkomnere, spiller en
stor Rolle, navnlig jo lavere Organismen staar. Men jo
mere udviklet den organiske Form er, des mere vil den
direkte Udviklingsproces, hvorved den ene Slægt overleverer
den næste Resultaterne af sine Kampe og Erfaringer,
nedlagte i Organisationen, afløse den indirekte.
Navnlig sker dette, naar et Nerve- og Muskelsystem har
udviklet sig, og dermed Evne til at variere Handlingerne
efter Omstændighederne. Naar Organismen er stærk nok,
vil den kunne læmpe sig; kun naar den er for svag, gaar
den under.

Der er her ifølge Spencer kun Valget imellem Udviklingstheorien
og Skabelsestheorien. Enten have Arterne udviklet
sig ifølge naturlige Love under naturlige Betingelsers
Vekselvirkning, eller ogsaa er hver Art skabt for sig efter en
fast Type. At lade hele Spørgsmaalet staa hen, lader sig
i Længden ikke gøre; og naar man i det Hele i sin Verdensanskuelse
støtte sig til naturlige Aarsager, saa kan Valget
mellem Udvikling og Skabelse ikke være vanskeligt. Skabelsen
er endog i sidste Instans utænkelig, kan end ikke
konstrueres for Fantasien: den udtrykker en Overgang fra
Intet til Noget, --- hvad der, som vi allerede have vist,
strider mod vor Tænknings og Bevidstheds Grundlov.
Man kan derfor sige, at Ingen i Virkeligheden tror paa
en Skabelse; for at kunne tro maa man dog i det Mindste
virkelig kunne forestille sig hvad man skal antage, men
% --- p. 119 ---
\setcounter{footnote}{0}%
\opage{119}det er her umuligt. Man tror, at man tror\footnote[1]{Smlgn, \emph{Feuerbachs} Epigram „Potenzirter Glaube“: Auf der höchsten Potenz steht jetzt in den Frommen der Glaube; Denn es wird ja sogar selber der Glaube geglaubt.}. Og hvor
mange enkelte Vanskeligheder vilde Skabelseshypothesen
ikke føre til! Skaberen maatte have fastholdt en Type,
selv hvor den var uhensigtsmæssig; Sælhunden f. Eks.
skulde være skabt til sit nuværende Opholdssted og dog
have Pattedyrslemmer, hvilket under Udviklingshypothesen
bliver ganske forklarligt, naar man tænker sig Sælen som
et Landdyr, der efterhaanden har akkommoderet sig efter
Livet i Havet. Det vilde endelig lige frem følge af Skabelseshypothesen,
at visse Arter vare bestemte til at tilføje andre
Smerte og Undergang. Udviklingshypothesen, der slet ikke
indlader sig paa den transcendente Aarsag ser i det organiske
Livs Smerte og Kamp et af de vigtigste Drivhjul,
der fører fremad til større Harmoni; den kan ganske
vist ikke --- saa lidt som nogensomhelst anden Lære
kan det --- sige os, hvorfor Smerte og Nød ere nødvendige
Drivhjul i Verdenslivet; men den fritager os for den
Tro, at de skulde være frembragte med Overlæg\footnote[2]{I Anledning af den skarpe Modsætning, Spencer hævder mellem Udvikling og Skabelse, har han havt en lille Polemik med Gladstone. Den berømte Statsmand og Lærde vilde nok gaa ind paa de almindelige Love og de naturlige Aarsager, men han bebrejdede Videnskaben og særlig Udviklingshypothesen, at den „indlod sig i Strid med Forsynet“; særlig udtalte han den Overbevisning, at Forsynet til behørig Tid lod de store Mænd opstaa, og nægtede, at dette kunde forklares alene af naturlige Love. Spencer sammenligner Gladstones Polemik mod naturlig Forklaring med Grækernes Harme over den fysiske Opfattelse af Himmellegemerne og deres Bevægelser. Særlig paa det sociale og personlige Livs Omraade gaa i vor Tid den naturlige og den teologiske Opfattelse Side om Side, uden at Modsigelsen mærkes. „Vi godtgjorde vor nationale Færdighed i at nære Forestillinger, der gensidig ophæve hinanden, da en upopulær Prinds pludselig vandt Popularitet ved at \opage{120}gennemgaa visse sygelige Forandringer i sit Blod, og da vi i Anledning af hans Helbredelse paa én Gang anerkendte Forsynets Hjælp og de naturlige Aarsagers Virken ved en Takkefest til Gud og en Baronettitel til Lægen.“ \emph{Spencer}: The study of sociology. p. 392---394; 423---426.}% the note runs over from p. 1119 to p. 1120 (joined by draft.py)
.

% --- p. 120 ---
\setcounter{footnote}{0}%
\opage{120}Udviklingshypothesen bringer tillige, som vi allerede
have set, alle vor Erkendelses Omraader i en storartet
Sammenhæng. Solsystemets, Jordklodens, Organismens,
Sprogets, Videnskabens og det menneskelige Samfunds
Historie gaar ind under den ene og samme almindelige
Lov. Under denne almindelige Lov gaa ogsaa Bevidsthedslivets
Fænomener ind. Psykologien er kun en Del af
Biologien, der igen er en udsondret Del af Geologien,
lige som denne atter er en speciel Del af Astronomien.
Vore Sondringer og Grænser ere subjektive og vilkaarlige
overfor den uendelige Sammenhæng i Universet. Alt er
i Kontinuitet. Saaledes kan man da heller ikke paavise
noget bestemt Punkt, hvor det bevidste Liv skiller sig
fra det ubevidste. Overgangen fra de fysiologiske til de
psykologiske Fænomener skér gennem den stigende Koncentration
af Nervesystemet, hvorved de spredte Paavirkninger
samles og bringes i Forbindelse med hverandre.
Herved bliver et højere Trin af den almindelige Livsproces
mulig: Indtrykkene kunne nu sammenlignes og
vurderes og en mere sammensat og mere adækvat Reaktion
bliver derved mulig fra Organismens Side. Bevidsthedslivet
udvikler sig Trin for Trin fra de laveste Former
til de højeste. Hos de laveste Dyr er kun dunkel Fornemmelse
og saa godt som ingen Kombination og Sammenligning
af de enkelte Fornemmelser mulig: Det er da
ogsaa kun meget ensidige og begrænsede Vekselvirkningsforhold
med Omverdenen, der her kunne finde Sted.
Man tænker sig f. Eks. Bevidsthedslivet hos en Indvoldsorm
eller en Østers. Jo rigere Organisationen bliver,
des nøjagtigere, mangesidigere og mere omfattende bliver
% --- p. 121 ---
\setcounter{footnote}{0}%
\opage{121}Vekselvirkningen. Den menneskelige Aand gennemtrænger
Jordkloden og Universet og har tillige i sig selv en rig
og mangfoldig Verden. Man sammenligne hin Indvoldsorms
Bevidsthed med Newtons, der udspreder Verdenslovene,
eller med Shakspeares, der gransker Sjælens dybeste
Drifter og Tanker. Og dog viser Dyre- og Menneskepsykologien
os en kontinuerlig Række af Mellemstadier
mellem disse Modsætninger.

Paa samme Maade forholder det sig indenfor selve
den menneskelige Aand. De forskellige saakaldte Evner
--- Instinkt, Fornuft, Iagttagelse, Begrebene, Erindring
etc. --- betegne kun forskellige Grader eller Former for
en og samme Proces. Fra den simpleste Kombination
af Fornemmelser eller Fremstillinger op til det mest abstrakte
og indviklede videnskabelige Ræsonnement er det
en og samme Grundvirksomhed, der ytrer sig: nemlig
Iagttagelse af Lighed og Ulighed mellem Indtryk eller
Forestillinger. Al Tænkning er en Klassifikation, og
Tankelivets Udvikling beror paa, hvor omfattende, sammensat
og nøjagtig denne Klassifikation er. Den højeste
Tænkningsform, den kvantitative eller mathematiske,
skylder sine Fortrin fremfor de andre Former den Omstændighed,
at den forudsætter Ensartethed hos sine
Genstande og kun undersøger Forholdet mellem deres
Størrelser; men den bestaar i sig selv lige saa vel i en
Iagttagelse af Lighed og Ulighed som den mest elementære
Sansning.

Den herskende Retning indenfor den engelske Filosofi,
som i vore Dage har havt en fremragende Repræsentant
i John Stuart Mill, var afgjort Empirisme, d. v. s.
den udledte al Erkendelse af Erfaringen og opfattede
Bevidstheden som en ren Tavle (tabula rasa), der efter
Haanden beskreven af Indtryk, saa igen kombineredes til
Forestillinger og Begreber. Imod denne Opfattelse erklærer
Spencer sig. Naar Aanden, siger han, forud for
al Erfaring er en tabula rasa, hvorfra kommer saa Evnen
% --- p. 122 ---
\setcounter{footnote}{0}%
\opage{122}til at organisere og bearbejde Erfaringerne? og hvorfra
komme de Forskelligheder, vi i denne Henseende finde
hos forskellige Racer og Individer? Hvorfor er saa i
det Hele et Nervesystem nødvendigt? Men det Vigtigste
er, at den gennemførte Empirisme vilde ophæve Erkendelsens
Sammenhæng og Nødvendighed. Naar det i det
Hele skal være muligt at lære af Erfaringen, saa maa
vi have et Princip, et Kriterium, efter hvilket vi kunne
maale den, og i Kraft af hvilket ethvert Fremskridt i
Erkendelse gøres. Dette sidste Kriterium finder Spencer
i Umuligheden af det Modsatte. Hvad vi ikke kunne
tænke, hvad der indeholder en Selvmodsigelse, kunne vi
heller ikke antage for sandt. Dette er den Prøve, hvortil
vi ved enhver Debat vises tilbage, og uden hvilken
Empirismen selv ikke kunde bevise en eneste af sine
Sætninger.

Uagtet Spencer saaledes giver Leibnitz og Kant Ret
i deres Polemik mod Empirismen, anser han det for lige
saa umuligt at betragte Erkendelsens Former som faststaaende,
indplantede uden Indflydelse fra den omgivende
Verden. Empirismens Fejl er, at den vil udlede Alt af
den \emph{individuelle} Erfaring. Dennes Utilstrækkelighed
viser os ud over Individet til Slægten. Det enkelte Individ
begynder ikke absolut forfra; det har i sin medfødte
Organisation Resultaterne af alle tidligere Generationers
Erfaringer, for saa vidt de ere gaaede over i
Strukturen. De universelle og gennemgaaende Former i
Tingene maa efterhaanden frembringe tilsvarende Former
i de Organismer, som i uendelige Tidsrækker ere underkastede
deres Indflydelse. Erkendelsens oprindelige Former
eller de ophobede og overleverede Strukturmodifikationer,
som frembringes ved Erfaring, ligge skjulte i ethvert
nyfødt Individ, røre sig paa en ubestemt Maade
ved de første individuelle Erfaringer og blive gradevis bestemtere
udprægede, naar saadanne individuelle Erfaringer
mangfoldiggøres.

% --- p. 123 ---
\setcounter{footnote}{0}%
\opage{123}Spencer kommer altsaa til det Resultat: a posteriori
for Slægten, men a priori for Individet; hvad der fra
først af er Erfaring, bliver gennem Organisationen til
Princip. Paa denne Maade søger han at komme ud over
Modsætningen mellem de to filosofiske Hovedretninger i
den nyere Tid, ud over Striden mellem Locke og Kant.

Dog regner han sig endnu selv til Lockes Skole,
fordi ifølge Udviklingshypothesen al Erkendelse stammer
fra Erfaringen, kun at det er Slægtens, ikke Individets
Erfaring. Hertil maa jeg bemærke, at Spencers Kritik
af Empirismen lige saa vel gælder, naar der er Tale om
Slægtens, som naar der er Tale om Individets Erfaring.
For at Slægten kan gøre Erfaringer, maa den have en
Organisation. Her er Tale om et væsentligt Forhold,
som er det samme paa ethvert Stadium. Desuden behandlede
baade Kant og Locke Spørgsmaalet udelukkende
fra Individets Synspunkt. Naar Spencer her giver Kant
Ret, staar han aabenbart med Hensyn til Problemet i
det Hele den tyske Filosof nærmere end han staar sine
engelske Forgængere. Kant lærer jo ogsaa udtrykkelig,
at al vor Erkendelse begynder med Erfaringen; kun negter
han, at den helt og holdent udspringer af Erfaringen\footnote[1]{I min udførligere Fremstilling af Spencers Lære i Skriftet „Den engelske Filosofi i vor Tid“ har jeg allerede gjort opmærksom paa, at Spencer trods sin egen Erklæring maa siges at staa Kant nærmere end Locke (p. 170 f.). Af det senere udkomne 3dje Bind af Spencers „Essays“ ser jeg, at \emph{Max Müller} har gjort en lignende Bemærkning. (p. 275: „In admitting that there is something in our mind, which is not the result of \emph{our own} à posteriori experience, Mr. Spencer is a thorough Kantian“).}.
Spencer indfører et helt nyt Synspunkt ved Erkendelsesproblemets
Behandling; derfor er det vanskeligt eller
umuligt at henføre ham til nogen af de Retninger, i
hvilke den tidligere Tænkning har spaltet sig under Debatten
om dette Problem. Jeg kan ikke se Andet end,
% --- p. 124 ---
\setcounter{footnote}{0}%
\opage{124}at han ved dette nye Synspunkt har bragt hele Spørgsmaalet
et betydeligt Skridt fremad\footnote[1]{Af et interessant Brev fra \emph{Stuart Mill} af 29de Januar 1872 til Fysiologen \emph{Carpenter}, som anføres i dennes „Mental Fysiology“ p. 486, ser man, at den rene Empirismes største Repræsentant i vor Tid mod Slutningen af sit Liv nærmede sig en lignende Opfattelse som Spencer.}.

Gælder det saaledes om Menneskets Erkendelse, at
den kun bliver forstaaelig, naar man gaar ud over det
enkelte Individ og ser dette i dets inderlige Sammenhæng
med hele Slægten, saa finder det dog i endnu højere
Grad Anvendelse, naar Talen er om Drift, Følelse og
Vilje. Drift og Følelse rører sig i Individet, førend dette
har gjort nogen Erfaring, der kunde frembringe dem;
navnlig gælder dette om den stærkeste af alle Drifter,
Kønsdriften. De sidste Motiver til den menneskelige
Handlen maa søges i de Grunddrifter, som Slægten har
erhvervet under sin Udviklingsgang. Ingen nok saa klar
Indsigt fremtvinger Handling uden ved at fremkalde en
tilstrækkelig stærk og varig Følelse. Det er altsaa ikke
sandt, at Ideerne regere Verden. Ideerne tjene kun til
Førere; det er den individuelle og nationale Karakter,
som bestemmer, hvorvidt Føreren følges. Ideerne blive
hovedsagelig til ved de Enkeltes Arbejde og kunne derfor
ikke direkte overvinde de dybe og mægtige Impulser,
som forplante sig ned gennem Slægten. Det følger tillige
af sig selv, at ingen Ideer kunne opstaa, som
staa i fuldstændig Modsætning til den herskende Karakter.
Naar højere Ideer have udviklet sig, bidrage de ganske
vist til videre Udvikling og Fremgang. Men den vigtigste
Faktor er og bliver dog selve den Tillæmpelse efter de
naturlige og sociale Vilkaar, som enhver Generation maa
forsøge. Karakteren udvikles kun ved Handling, og
Handling fremkaldes ved de virkelige Forhold. Den sande
Opdragelse --- indenfor Familjen som indenfor Nationen
% --- p. 125 ---
\setcounter{footnote}{0}%
\opage{125}i det Hele --- maa derfor bestaa i, at Individerne holdes
i Vekselvirkning med de faktiske Betingelser.

Men denne Indsigt er ifølge Spencer saare langt fra
at være almindelig. Af alle de Fordomme, som hindre
sund Tænkning paa Opdragelsens og Samfundslivets Omraade,
er ingen mere udbredt og mere skadelig end den,
at Forældrene og Regeringerne kunne udrette noget Positivt.
Ved direkte at gribe ind hindre de kun den naturlige
Vekselvirkning mellem Individerne og Omgivelserne.
De have kun at sørge for, at de nødvendige Betingelser
ere til Stede for Individets Bestaaen og Udvikling. Den
Forudsætning, hvorfra man i Almindelighed gaar ud, er,
at Dyderne findes hos de Regerende, Lasterne hos de
Regerede. Man ser ene hen til den \emph{ordnende} Kraft
og glemmer, at det Væsentlige dog altid maa være selve
det, som \emph{ordnes}.

De moderne Reformatorer, som mene at kunne frelse
Menneskeheden ved deres Undervisningssystemer eller
Statsforfatninger, ere hildede i ikke mindre Vildfarelse
end de Konservative og Reaktionære; begge Parter have
det til Fælles, at de vilde bringe de Enkelte under Formynderskab.
Det er en Slags Fetischdyrkelse af Regeringsmagten,
der har udviklet sig, og som Radikale og
Toryer ere lige befængte med. Menneskene kunne dog
ikke frembringe nye Kræfter; de maa virke med de
Midler, Naturen giver. Samfundets Karakter beror paa
Karakteren af de Individer, ved hvis Forening det bestaar.
Men denne Karakter forandres ikke paa én Gang
ved at indøves i visse intellektuelle Færdigheder\footnote[1]{„Uvidenhed og Forbrydelse ere ikke Aarsag og Virkning, men Virkninger af en og samme Aarsag. At være \emph{helt} uden Kundskaber vil sige at have levet mellem Folk, som have de stærkeste Impulser til Forbrydelse; at have modtaget \emph{nogen} Undervisning vil sige at høre til en Klasse, som er mindre stærke Fristelser underkastet; at være \emph{vel} undervist vil sige at have levet under Forhold, hvor de sædvanlige Motiver til \opage{126}Forbrydelse næsten ikke røre sig.“ \emph{Spencer}: Social Statics. p. 379 f.}% the note runs over from p. 1125 to p. 1126 (joined by draft.py)
 eller
% --- p. 126 ---
\setcounter{footnote}{0}%
\opage{126}ved at bringes i et vist udenfra hentet System. Lige
saa lidt som man kan undgaa Tyngdekraftens Indflydelse
ved at anbringe sine Lemmer i en vis bestemt Stilling,
lige saa lidt kunne Menneskene haabe at undgaa den
Slethed, der tynger dem ned, ved at ordne sig efter dette
eller hint politiske Mønster. Man overser bestandig de
skjulte, i Dybet virkende Kræfter og stirrer paa de Former
og Faktorer, der bevæge sig paa Overfladen.

Urmennesket maatte have vilde og egoistiske Instinkter
for at kunne bestaa i den haarde Kamp imod Naturen,
Dyrene og sine Lige. Da et Samfund begyndte at danne
sig, --- et vigtigt Middel til at føre denne samme Kamp,
--- blev der paalagt Individerne visse Baand, der vare
nødvendige naar Samfundet i det Hele skulde bestaa.
Det Samfund og de Individer kom videst og bestod bedst
i Kampen, hos hvem Tillæmpelsen efter Samfundslivets
Fordringer skred videst frem. De egoistiske Instinkter
bleve derfor ikke længe eneraadende; de sympatetiske
Drifter, der allerede fra først af have saa dyb en Grund
i Kønsforholdet og i Forholdet mellem Forældre og Børn,
maatte mere og mere tiltage i Styrke og Omfang. Den
Ene agter den Andens Ret og arbejder maaske endog
uegennyttig for hans Vel. Egoismen gaar saaledes over
i Altruisme, for at bruge Comtes træffende Udtryk, og
den sociale Udviklings Ideal, der ogsaa indeholder det
ethiske Ideal, er en Tilstand, hvor Alle virke sammen,
Enhver efter sin Individualitet og med saa megen Frihed,
som kan bestaa med alle Andres Frihed.

Men dette er en Tilstand, som kun den naturlige
Udvikling kan føre til. Regeringen skal kun sørge for
Sikkerhed og Retfærdighed, Grundbetingelserne for, at
Udviklingen kan gaa for sig, Vistnok er i Samfundet
som i enhver Organisme en vis Organisation, en vis
% --- p. 127 ---
\setcounter{footnote}{0}%
\opage{127}Struktur nødvendig. Men jo mindre denne fremgaar af
selve Elementernes Væsen og Vekselvirkning, jo mere
kunstig og ovenfra indsat den er, des større Hindring vil
den fremadskridende Udvikling have at overvinde ved
den. Sprogets Udvikling og den økonomiske Udvikling,
de frembringende, ombyttende og fordelende Processer,
ere Eksempler paa Resultater, der ere naaede ad naturlig
Vej, hvor hver Enkelt ved at arbejde under de givne
Betingelser og efter de naturlige Formaal frembragte
Resultater, der laa langt ud over den individuelle Synskreds
baade i Omfang og i Betydning. Efterhaanden,
som de egoistiske Instinkter trænges tilbage, ville ogsaa
uegennyttige Bestræbelser hos Enkelte eller frie Forbindelser
af Enkelte, lægge deres Lod i Vægtskaalen og udrette,
hvad hine ikke formaaede. Staten, som ikke engang
kan overkomme sin Hovedopgave, en upartisk og energisk
Hævdelse af Ret og Sikkerhed, kan langt mindre udrette
noget paa andre Omraader og vil i hvert Tilfælde gøre mere
Skade end Gavn. Menneskeheden er i en Overgangstilstand;
den skal aflægge sin oprindelige Røvernatur og
blive skikket til fuldkomment Samfundsliv. Og denne
Overgangstilstand maa nødvendigvis være ulykkelig: al
Ulykke bestaar i Uoverensstemmelse mellem den indre
Organisation og de ydre Betingelser. Men Lidelserne
maa gennemgaas. De kunne lindres ved Sympati, ved
den Kærlighed, der hjælper Andre til at hjælpe sig
selv; men et vist Maal af Lidelse er uundgaaeligt. Voldsomt
at gribe ind vilde netop forøge og forlænge Lidelsen.
Der er Larver, som leve i Vandet og aande ved Gæller,
men hos hvem der allerede her udvikles Lunger til Brug
for en senere Periode i Dyrets Liv. Førend Lungerne
ere fuldt udviklede, kan Dyret ikke undvære Gællerne;
man vilde dræbe det, dersom man pludselig vilde gøre
Ende paa den Inkonsekvens at have Lunger og Gæller
paa én Gang. Og saaledes gaar det ogsaa i Samfundslivets
Udvikling; vi kunne forhale denne ved direkte at
% --- p. 128 ---
\setcounter{footnote}{0}%
\opage{128}
gribe ind, men vi kunne ikke fremskynde den paa anden
Maade end ved at hævde de naturlige Betingelsers Indflydelse.
Og selv i de Lidelser, der herved nødvendigvis
indtræder, vil den, der fæster sit Blik paa Udviklingen
og dens Maal i det Hele, erkende en kæmpende Kærlighed.

Spencers ethiske og sociologiske Anskuelser foreligge
endnu ikke i udarbejdet Form. Af hans filosofiske System,
der skal udgøre 10 Bind, staar netop Halvdelen
tilbage, som skal indeholde Sociologiens og Etikens Principer\footnote[1]{De fem første Bind ere: First Principles, Principles of Biology (2 Bd.) og Principles af Psychology (2 Bd.).}.
Men den almindelige Karakter og Retning fremlyser
tydelig af hvad vi kende. Det Første, jeg vil fremhæve,
er, at Spencer lægger an paa en aldeles objektiv,
streng kausal Opfattelse af det menneskelige Samfundsliv.
Han er ikke, som Stuart Mill er i alle sine sociale og
politiske Afhandlinger, tillige Reformator. Dette maa
man vel have for Øje. Man har beskyldt ham for Fatalisme,
fordi han lægger saa megen Vægt paa, at den
menneskelige Udvikling er en Naturudvikling, og for nylig
har en Discipel af Stuart Mill anført imod ham følgende
Ord af denne: „Politiske Institutioner ere Menneskenes
Værk, skylde den menneskelige Vilje deres Oprindelse
og hele deres Eksistens. Menneskene vaagne ikke op
en skøn Morgen og finde dem fremvoksede\footnote[2]{\emph{Cairnes}: Fortnightly Review. January 1875.}.“ Men Spørgsmaalet
er jo ifølge Spencer netop, om Menneskene ikke
i selve deres Villen og Virken adlyde en Naturlov, de
ikke ere sig bevidste. Han vil atter her citere Shakspeares
Ord, at den Kunst, der siges at hjælpe paa Naturen,
har sin Grund i en Kunst, Naturen selv har frembragt.
Og man vil nødes til at indrømme dette, med
mindre man vil hævde et Begreb om Viljens Frihed, som
sprænger al Lovmæssighed og Sammenhæng paa det
psykologiske og sociale Omraade.

% --- p. 129 ---
\setcounter{footnote}{0}%
\opage{129}Spencers Lære vil heller ikke trænge den bevidste
Stræben og Virken tilbage; det er kun Vejen og Retningen
for denne Stræben, Uenigheden drejer sig om. Efter
hans Overbevisning kan den sande og fulde Udvikling af
Menneskets personlige Kræfter kun finde Sted, naar man
ikke hindrer den Enkelte i at komme Ansigt til Ansigt
med de virkelige Forhold og optage Kampen for Tilværelsen
under de reale og faktiske Vilkaar. Statsmagtens
Opgave bestaar da for ham, lige som for Wilhelm
v. Humboldt, i en Emancipation af den individuelle
Frihed. Op paa Grund af de Udviklingsbetingelser, der
kræves for den sociale Livsproces som for enhver Livsproces,
er det efter hans Overbevisning en Illusion at
vente betydelige Resultater i korte Tisrum\footnote[1]{Den berømte Embryolog \emph{Ernst} v. \emph{Baer}, af hvem Spencer har lært saa meget, er kommen til lignende Resultat ved at sammenligne den sociale Udvikling med den organiske. „Es prägte sich mir tief die Ueberzeugung ein, dass die socialen Entwickelungen nicht anders vor sich gehen können als die organischen, nämlich allmählich, und nur dann gesunde Gestalt gewinnen, wenn die Bedingungen dazu vorher vollständig gegeben waren und längere Zeit bestanden hätten \dots{} Die Männer, welcke vorwärts wollen, --- und rückwärts zu wollen ist überhaupt gegen die Natur, --- würden gut thun, ihre Gedanken recht viel an die Beachtung der organischen Entwickelungen zu gewöhnen.“ Nachricht über Leben und Schriften K. E. v. Baer’s, mitgetheilt von ihm selbst. St. Petersburg 1865. p. 52. --- I sit sidste Skrift „The study of sociology“ fordrer Spencer som nødvendig Forberedelse til Beskæftigelse med Samfundsvidenskaben Studium af de eksakte Videnskaber, af Biologi og Psykologi.}. % PRINTED AS IS (Tisrum: for Tidsrum; reader b1-04-01)

Men det kan ikke negtes, at Spencer overser den
positive Betydning, som Statsmagtens Indgriben kan have
for Opnaaelsen af det Formaal, hvorefter den sociale
Udvikling stræber. Vel vilde han næppe give Buckle Ret
i, „at det vilde være at drive Spot med al sund Fornuft,
dersom man tilskrev Lovgivningen nogensomhelst Andel
% --- p. 130 ---
\setcounter{footnote}{0}%
\opage{130}i Fremskridtet eller af fremtidige Lovgivere ventede
nogen anden Velgerning end Tilbagekaldelsen af deres
Forgængeres Værk;“ --- vel udtaler han (i Fortalen til
anden Udgave af „Social Statics“), at han nu mere end
før vilde betone de politiske Institutioners relative
Værdi; --- men man savner dog hos ham den klare
Erkendelse af Samfundets opdragende Betydning. „Der
er,“ som jeg et andet Sted\footnote[1]{Den engelske Filosofi i vor Tid, p. 65.} har udtrykt denne, „endnu
en lang Vej at gaa, førend den i Sandhed naturlige Udvikling
kan indtræde. Staten, ved hvis Hjælp Mennesket
er ført ud af den dyriske Naturtilstand, kan ikke slippe
ham, førend den har ført ham til den humane Naturtilstand.“
Det er herpaa, at de socialistiske Ideers relative
Berettigelse grunder sig. Ogsaa her ligger Løsningen
af Modsætningerne i Udviklingens Ide; Modsætningen
mellem Individualisme og Socialisme, der efter
Stuart Mills Udsagn\footnote[2]{John Stuart Mills Levnet fortalt af ham selv. Oversat af V. Arntzen, p. 259.} er diametral, kan kun overvindes
ved, at begge Retninger opfattes som gensidig betingende
hinanden. Men fra en saadan Løsning ere vi endnu langt
borte. Spencers Sociologi, der overvejende bevæger sig
paa det teoretiske Omraade, fører til en gavnlig Skepsis
overfor alle foreløbige Løsningsforsøg, samtidig med at
den energisk peger mod de dybest virkende Kræfter og
det ene, væsentlige Maal.
\signature{\textbf{Harald Høffding.}}

\end{document}
